% ==============================================================================
% APÊNDICES - CÓDIGOS-FONTE DO PLANO DE DADOS SPOLKA EM eBPF/XDP
% ==============================================================================
% Arquivo formatado para ser incluído diretamente no main.tex da classe ifes8:
%   \printbibliography
%   \include{tex/apendices/apendices_ebpf}

\begin{apendicesenv}

\chapter{CÓDIGO-FONTE DO ENCAPSULAMENTO NA BORDA (\texttt{spolka\_edge\_encap.c})}
\label{apendice:encap}

Este apêndice apresenta o programa eBPF anexado ao gancho \texttt{TC egress} da interface \texttt{eth1} dos nós de borda ($H_1$--$H_3$ e $S_1$--$S_3$). O código realiza a classificação dos pacotes IPv4 mediante consulta de prefixo mais longo (\texttt{BPF\_MAP\_TYPE\_LPM\_TRIE}) sobre a chave composta pelo endereço IPv4 de destino e pelo campo DSCP, inserindo os 8~bytes do identificador \routeID\ no cabeçalho do quadro por meio da função auxiliar \texttt{bpf\_skb\_change\_head()}.

\begin{lstlisting}[language=C, caption={Programa eBPF de classificação e encapsulamento na borda (\texttt{spolka\_edge\_encap.c}).}, label={lst:spolka_edge_encap}]
#include <linux/bpf.h>

#include <bpf/bpf_endian.h>
#include <bpf/bpf_helpers.h>

#include <linux/if_ether.h>
#include <linux/ip.h>
#include <linux/tcp.h>
#include <linux/pkt_cls.h>

#ifndef IPPROTO_TCP
#define IPPROTO_TCP 6
#endif

#define TYPE_SRCROUTING 0x1234
#define ROUTE_ID_SIZE 8

// Chave da tabela_encap: (IP de destino, DSCP). O IP vem ANTES do TOS de
// proposito. O LPM trie casa um prefixo contiguo a partir do bit 0, entao:
//
//   prefixlen 40 -> IP + classe   (routeID especifico daquela classe)
//   prefixlen 32 -> IP            (default: vale para qualquer classe)
//
// Como o casamento mais longo vence, uma unica busca resolve os dois casos:
// a classe com routeID proprio pega a entrada de 40 bits, e qualquer outra
// cai no default de 32 bits. E o que permite o slicing por classe do SPolKA
// sem precisar instalar as 7 classes em todo destino.
//
// Os 3 bytes de padding no fim da struct ficam fora do prefixo (bits 40..63)
// e portanto nao participam do casamento.
struct polka_lpm_key {
  __u32 prefixlen;
  __u32 ipv4_addr;
  __u8 tos;
};

#define PREFIX_IP_BITS 32
#define PREFIX_IP_TOS_BITS 40

// Mascara os 2 bits de ECN: o TCP os altera em voo, e so o DSCP identifica a
// classe de servico.
#define DSCP_MASK 0xFC

struct {
  __uint(type, BPF_MAP_TYPE_LPM_TRIE);
  __uint(max_entries, 1024);
  __type(key, struct polka_lpm_key);
  __type(value, __u64);
  __uint(map_flags, BPF_F_NO_PREALLOC);
} tabela_encap SEC(".maps");

SEC("classifier")
int polka_edge_encap(struct __sk_buff *skb) {
  void *data_end = (void *)(long)skb->data_end;
  void *data = (void *)(long)skb->data;

  struct ethhdr *eth = data;
  if ((void *)(eth + 1) > data_end)
    return TC_ACT_OK;

  if (eth->h_proto != bpf_htons(ETH_P_IP)) {
    return TC_ACT_OK;
  }

  struct iphdr *iph = (void *)(eth + 1);
  if ((void *)(iph + 1) > data_end)
    return TC_ACT_OK;

  struct polka_lpm_key key = {
      .prefixlen = PREFIX_IP_TOS_BITS,
      .ipv4_addr = iph->daddr,
      .tos = iph->tos & DSCP_MASK,
  };

  __u8 *route_id_encontrado = bpf_map_lookup_elem(&tabela_encap, &key);
  if (!route_id_encontrado) {
    bpf_printk("encap: no route ID found for dest %pI4 tos 0x%x\n", &iph->daddr,
               key.tos);
    return TC_ACT_OK;
  }

  // --- EXECUCAO DO PUSH DO CABECALHO POLKA ---

  // 1. Salva o cabecalho Ethernet original na pilha
  struct ethhdr eth_backup;
  __builtin_memcpy(&eth_backup, eth, sizeof(struct ethhdr));
  eth_backup.h_proto = bpf_htons(TYPE_SRCROUTING);

  __u8 route_id_backup[ROUTE_ID_SIZE];
  __builtin_memcpy(route_id_backup, route_id_encontrado, ROUTE_ID_SIZE);

  // 2. Prepend de 8 bytes na cabeca do pacote. bpf_skb_change_head so desloca
  // o ponteiro de dados, sem tocar em metadados de L3/L4 (checksum, GSO) --
  // diferente de bpf_skb_adjust_room, que corrompe esses metadados e quebra
  // TCP (segmentos presos, sem avanco de cwnd).
  if (bpf_skb_change_head(skb, ROUTE_ID_SIZE, 0)) {
    bpf_printk("encap: bpf_skb_change_head failed\n");
    return TC_ACT_OK;
  }

  data = (void *)(long)skb->data;
  data_end = (void *)(long)skb->data_end;

  struct ethhdr *new_eth = data;
  if ((void *)(new_eth + 1) > data_end) {
    bpf_printk("encap: new_eth out of bounds\n");
    return TC_ACT_SHOT;
  }

  // 3. Escreve o novo Ethernet header (offset 0)
  __builtin_memcpy(new_eth, &eth_backup, sizeof(struct ethhdr));

  // 4. Escreve o RouteID logo apos o novo MAC (offset 14)
  __u8 *polka_route_field = (__u8 *)(new_eth + 1);
  if ((void *)(polka_route_field + ROUTE_ID_SIZE) > data_end) {
    bpf_printk("encap: polka_route_field out of bounds\n");
    return TC_ACT_SHOT;
  }

  __builtin_memcpy(polka_route_field, route_id_backup, ROUTE_ID_SIZE);

  return TC_ACT_OK;
}

char _license[] SEC("license") = "GPL";
\end{lstlisting}
\fonte{Elaborado pelo autor (2026).}

% ==============================================================================
\chapter{CÓDIGO-FONTE DO ENCAMINHAMENTO NO NÚCLEO (\texttt{spolka\_core.c})}
\label{apendice:core}

Este apêndice apresenta o programa eBPF anexado ao gancho \texttt{TC ingress} (\texttt{clsact}) das interfaces de dados dos roteadores de núcleo ($C_1$ a $C_8$). O programa executa a divisão polinomial em $GF(2)$ (\texttt{poly\_mod}) sobre o \routeID\ utilizando o polinômio irredutível local armazenado em \texttt{mapa\_polinomio}, extrai simultaneamente a porta lógica de saída e o identificador da fila (\texttt{queue\_id}), registra a fila em \texttt{skb->priority} e redireciona o pacote via \texttt{bpf\_redirect()}.

\begin{lstlisting}[language=C, caption={Programa eBPF de decodificação em $GF(2)$ e encaminhamento sem estado no núcleo (\texttt{spolka\_core.c}).}, label={lst:spolka_core}]
#include <linux/bpf.h>

#include <bpf/bpf_endian.h>
#include <bpf/bpf_helpers.h>

#include <linux/if_ether.h>
#include <linux/pkt_cls.h>

#define TYPE_SRCROUTING 0x1234

// Estrutura PolKA compacta
struct polka_t_top {
  unsigned char h_dest[ETH_ALEN];
  unsigned char h_source[ETH_ALEN];
  __be16 etherType;
  __u8 routeId[8];
} __attribute__((packed));

struct {
  __uint(type, BPF_MAP_TYPE_ARRAY);
  __uint(max_entries, 1);
  __uint(key_size, sizeof(__u32));
  __uint(value_size, sizeof(__u16));
} mapa_polinomio SEC(".maps");

struct {
  __uint(type, BPF_MAP_TYPE_ARRAY);
  __uint(max_entries, 16);
  __uint(key_size, sizeof(__u32));
  __uint(value_size, sizeof(__u32));
} mapa_portas SEC(".maps");

/*
 * GF(2) Polynomial Modulo
 * Calcula: div % mod
 * Implementa divisao polinomial binaria (CRC-like) inteiramente sobre
 * registradores de 32 bits, sem acesso a memoria.
 */
static __always_inline __u32 poly_mod(__u32 div, __u32 mod) {
  if (mod == 0) return 0;

  // Encontra o bit mais significativo (MSB) do modulo
  int mod_msb = 0;
#pragma unroll
  for (int i = 0; i < 32; i++) {
    if ((mod >> i) & 1) {
      mod_msb = i;
    }
  }

  // Processa bits do MSB para baixo
#pragma unroll
  for (int i = 31; i >= 0; i--) {
    if (i >= mod_msb) {
      if ((div >> i) & 1) {
        __u32 shift = i - mod_msb;
        __u32 shifted_mod = mod << shift;
        div ^= shifted_mod;
      }
    }
  }
  return div;
}

SEC("classifier")
int polka_core_router(struct __sk_buff *skb) {
  void *data_end = (void *)(long)skb->data_end;
  void *data = (void *)(long)skb->data;

  struct polka_t_top *polka = data;
  if ((void *)(polka + 1) > data_end)
    return TC_ACT_OK;
  if (polka->etherType != bpf_htons(TYPE_SRCROUTING))
    return TC_ACT_OK;

  // RouteID de 64 bits: bytes superiores carregam o dividendo e os 2 bytes
  // inferiores carregam o termo independente (dif)
  __u16 dif = (__u16)((polka->routeId[6] << 8) | polka->routeId[7]);

  __u32 route_val = ((__u32)polka->routeId[0] << 24) |
                    ((__u32)polka->routeId[1] << 16) |
                    ((__u32)polka->routeId[2] << 8)  |
                    ((__u32)polka->routeId[3]);

  __u32 key = 0;
  __u16 *polinomio_no = bpf_map_lookup_elem(&mapa_polinomio, &key);
  __u32 poly = polinomio_no ? (__u32)*polinomio_no : 0xA001;

  __u32 nresult = poly_mod(route_val, poly);
  __u16 nport_raw = (__u16)nresult ^ dif;

  // --- EXTRACAO SIMULTANEA DE PORTA E FILA (BIT-SLICING) ---
  __u16 queue_id = nport_raw & 0x0007; // Extrai os 3 bits de Fila
  __u16 logical_port = nport_raw >> 3; // Extrai a Porta Logica

  __u32 key_port = logical_port;
  __u32 *ifindex_ptr = bpf_map_lookup_elem(&mapa_portas, &key_port);
  __u32 ifindex_out = ifindex_ptr ? *ifindex_ptr : 0;

  if (ifindex_out == 0)
    return TC_ACT_SHOT;

  // Informa a disciplina de fila (qdisc) a qual fatia este pacote pertence
  if (queue_id == 0) {
      queue_id = 1;
  }
  skb->priority = 0x00010000 | queue_id;

  // Redireciona diretamente para o Egress da interface de saida
  return bpf_redirect(ifindex_out, 0);
}

char _license[] SEC("license") = "GPL";
\end{lstlisting}
\fonte{Elaborado pelo autor (2026).}

% ==============================================================================
\chapter{CÓDIGO-FONTE DO DESENCAPSULAMENTO NA BORDA (\texttt{spolka\_edge\_decap.c})}
\label{apendice:decap}

Este apêndice apresenta o programa XDP nativo anexado à interface \texttt{eth1} dos nós de borda de egresso ($H_1$--$H_3$ e $S_1$--$S_3$). Ao interceptar quadros com \textit{EtherType} \texttt{0x1234}, o programa aciona a primitiva \texttt{bpf\_xdp\_adjust\_head()} para remover os 8~bytes do cabeçalho \nome, restaura o \textit{EtherType} para IPv4 (\texttt{0x0800}) e entrega o pacote original à pilha de rede do sistema operacional (\texttt{XDP\_PASS}).

\begin{lstlisting}[language=C, caption={Programa XDP nativo de remoção do rótulo na borda (\texttt{spolka\_edge\_decap.c}).}, label={lst:spolka_edge_decap}]
#include <linux/bpf.h>
#include <bpf/bpf_endian.h>
#include <bpf/bpf_helpers.h>
#include <linux/if_ether.h>
#include <linux/ip.h>

#define TYPE_SRCROUTING 0x1234
#define ROUTE_ID_SIZE 8

SEC("xdp")
int polka_edge_decap(struct xdp_md *ctx) {
  void *data_end = (void *)(long)ctx->data_end;
  void *data = (void *)(long)ctx->data;

  struct ethhdr *eth = data;
  if ((void *)(eth + 1) > data_end)
    return XDP_PASS;

  if (eth->h_proto != bpf_htons(TYPE_SRCROUTING)) {
    return XDP_PASS;
  }

  // --- EXECUCAO DO POP (REMOCAO DO CABECALHO) ---

  // 1. Guarda o cabecalho Ethernet original na memoria
  struct ethhdr eth_backup;
  __builtin_memcpy(&eth_backup, eth, sizeof(struct ethhdr));

  // 2. Encolhe o pacote em 8 bytes na cabeca (bpf_xdp_adjust_head)
  int err = bpf_xdp_adjust_head(ctx, ROUTE_ID_SIZE);
  if (err) {
    bpf_printk("decap: bpf_xdp_adjust_head failed: %d\n", err);
    return XDP_DROP;
  }

  // 3. Recalcula ponteiros (obrigatorio para o Verificador eBPF)
  data = (void *)(long)ctx->data;
  data_end = (void *)(long)ctx->data_end;

  struct ethhdr *new_eth = data;
  if ((void *)(new_eth + 1) > data_end) {
    return XDP_DROP;
  }

  // 4. Altera a copia guardada para o EtherType IPv4
  eth_backup.h_proto = bpf_htons(ETH_P_IP);

  // 5. Cola o cabecalho Ethernet guardado na nova posicao
  __builtin_memcpy(new_eth, &eth_backup, sizeof(struct ethhdr));

  return XDP_PASS;
}

char _license[] SEC("license") = "GPL";
\end{lstlisting}
\fonte{Elaborado pelo autor (2026).}

\end{apendicesenv}

